% ============================================================================
\section{Every two-level Hamiltonian is a spin}
\label{sec:formalism}
% ============================================================================

Chemistry supplies two-level systems in abundance. A bond in a minimal basis has two
atomic orbitals; a proton in an NMR magnet has two spin projections; a donor--acceptor
complex has the electron on one site or the other; ammonia has its umbrella up or down;
a chromophore in a weak laser field has a ground and a first excited state. In each case
the relevant Hilbert space is two dimensional, and in each case --- as we now show ---
the Hamiltonian is a linear combination of the identity and the three Pauli matrices.

\subsection{The Pauli matrices are a complete basis}

Let $\ket{\alpha}$ and $\ket{\beta}$ be an orthonormal pair spanning the space. Which
two physical states these are will change from section to section; the algebra will not.
Any Hamiltonian is a $2\times2$ Hermitian matrix,
%
\begin{equation}
\Hop =
\begin{bmatrix}
H_{\alpha\alpha} & H_{\alpha\beta}\\
H_{\alpha\beta}^{*} & H_{\beta\beta}
\end{bmatrix},
\qquad H_{\alpha\alpha}, H_{\beta\beta}\in\mathbb{R},
\label{eq:generalH}
\end{equation}
%
which carries four real degrees of freedom: the two diagonal entries and the real and
imaginary parts of $H_{\alpha\beta}$. The set $\{\Iop,\sxo,\syo,\szo\}$ carries four
real degrees of freedom as well, and it is straightforward to check that these four
matrices are linearly independent. They therefore span the space of Hermitian
$2\times2$ matrices, and we may always write

\begin{keyresult}[The Pauli decomposition]
\vspace{-0.35cm}
\begin{equation}
\Hop \;=\; \varepsilon_0\,\Iop \;+\; \half\!\left(d_x\sxo + d_y\syo + d_z\szo\right)
    \;\equiv\; \varepsilon_0\,\Iop + \half\,\dvec\cdot\pauliv ,
\label{eq:pauli-decomp}
\end{equation}
with $\varepsilon_0$ and the three components of $\dvec$ all real.
\end{keyresult}

\noindent
Writing \cref{eq:pauli-decomp} out in components,
%
\begin{equation}
\Hop =
\begin{bmatrix}
\varepsilon_0 + \tfrac{d_z}{2} & \tfrac{d_x - \ii d_y}{2}\\[4pt]
\tfrac{d_x + \ii d_y}{2} & \varepsilon_0 - \tfrac{d_z}{2}
\end{bmatrix},
\label{eq:pauli-components}
\end{equation}
%
so that each piece can be read off and given a name.

\begin{center}
\small
\begin{tabular}{@{}lll@{}}
\toprule
Component & Definition & Physical meaning \\
\midrule
$\varepsilon_0$ & $\half(H_{\alpha\alpha}+H_{\beta\beta})$
  & mean energy; shifts everything, changes nothing \\
$d_z$ & $H_{\alpha\alpha}-H_{\beta\beta}$
  & \textbf{asymmetry}: the energy gap between the basis states \\
$d_x$ & $2\,\mathrm{Re}\,H_{\alpha\beta}$
  & \textbf{coupling}: the real part of the mixing element \\
$d_y$ & $-2\,\mathrm{Im}\,H_{\alpha\beta}$
  & \textbf{coupling}: the imaginary part \\
\bottomrule
\end{tabular}
\end{center}

That table is the thesis of these notes. Everything that follows is an exercise in
identifying, for some piece of chemistry, what plays the role of the asymmetry and what
plays the role of the coupling.

\subsection{Extracting the components}

The Pauli matrices satisfy the two trace identities
%
\begin{equation}
\Tr\!\left(\sxo\right)=\Tr\!\left(\syo\right)=\Tr\!\left(\szo\right)=0,
\qquad
\Tr\!\left(\hat\sigma_i\hat\sigma_j\right) = 2\,\delta_{ij},
\label{eq:traces}
\end{equation}
%
which follow immediately from the product rule
$\hat\sigma_i\hat\sigma_j = \delta_{ij}\Iop + \ii\,\epsilon_{ijk}\hat\sigma_k$
(derived in \cref{app:pauli}). Multiplying \cref{eq:pauli-decomp} by $\hat\sigma_i$ and
tracing therefore isolates a single component:

\begin{equation}
\boxed{\;\varepsilon_0 = \half\Tr\!\left(\Hop\right),
\qquad
d_i = \Tr\!\left(\hat\sigma_i\,\Hop\right).\;}
\label{eq:extract}
\end{equation}

\begin{nbbox}[In the notebook]
\Cref{eq:extract} is the function \texttt{pauli\_decomposition(H)}, and
\cref{eq:pauli-decomp} is \texttt{build\_two\_level(e0, dz, dx, dy)}. The notebook
verifies that the two are exact inverses of one another on 500 randomly generated
Hermitian matrices --- a much stronger check than any single worked example, because
it tests completeness of the basis rather than one lucky case.
\end{nbbox}

\subsection{A first example: the spin operators themselves}

As a sanity check, take the operators from Computational Set 1. Since
$\Szo = \tfrac{\hbar}{2}\szo$, applying \cref{eq:extract} gives
$\varepsilon_0=0$ and $\dvec = (0,0,\hbar)$. Similarly $\Sxo$ gives
$\dvec=(\hbar,0,0)$ and $\Syo$ gives $\dvec=(0,\hbar,0)$. The three spin operators are
simply the three coordinate axes of the $\dvec$ space, which is why that space is worth
taking seriously as a geometric object --- the subject of the next section.
